% Introduction (BGU \S 10.7.2.4).
% V6 readability pass: shorter sentences, plain-language framing woven in,
% structure unchanged (Motivation / Problem / Alternatives / Contributions /
% Roadmap).
\section{Introduction}
\label{sec:intro}

\subsection{Motivation}
Terahertz (THz) imaging is increasingly used in security screening, biomedical
diagnostics, and through-fabric inspection, because radiation in this band
passes through clothing, paper, plastics, and other materials that block
visible light. The same physics that makes THz attractive also limits it:
diffraction caps the spatial resolution near the wavelength of the source,
the detector integrates each measurement over a comparatively large area, and
thermal background dominates the noise floor. The practical result is easy to
state: a THz frame is simultaneously very low resolution, noisy, and
colorless.

Standard deep-vision classifiers are built for the opposite kind of picture.
Almost all of them inherit ImageNet pretraining and rely on fine textures,
sharp edges, and saturated color. When those cues are removed, accuracy
collapses --- and, importantly, the collapse is opaque. Before this project it
was not obvious how much of the loss comes from the low resolution alone, how
much from noise or from the missing color, and whether anything on the
training side (regularization, augmentation, or simply choosing a different
architecture) can win part of the loss back. Answering those questions in a
controlled, reproducible way is the purpose of this work.

\subsection{Problem statement}
This project quantifies the accuracy degradation of three representative deep
architectures under controlled synthetic THz-like degradation, decomposes the
loss into per-axis contributions, audits the reproducibility of the
measurements across random seeds, and documents which training-side
interventions recover lost accuracy. The three architectures span the
families that anchor the image-classification literature: a residual
convolutional network (ResNet50 [1]), a densely connected convolutional
network with explicit feature reuse (DenseNet121 [2]), and a vision
transformer with foveal attention (TransNeXt-tiny [3]).

The success criteria carried over from the preliminary report are: at least
80 percent top-1 accuracy at the moderate-degradation operating point, an
identifiable architectural winner that beats the baseline by a statistically
meaningful margin, and inference latency compatible with sensor-side
deployment. These targets are revisited in Section~\ref{sec:conclusions}
against the actual measurements --- including the one target the measurements
show to be physically out of reach, which is itself an informative result.

\subsection{Alternative approaches considered}
Three alternative families were reviewed during the preliminary design and
rejected in favor of training directly on degraded images.

\textbf{(i) Denoise, then classify.}\quad
Restore the image first, then hand it to a clean-image classifier. The appeal
is modularity: the pretrained classifier is reused as-is. The problem is that
restoration models need matched clean/degraded training pairs, which a real
THz sensor never produces, and restoration networks trained on mismatched
degradations tend to hallucinate detail [4], [5] --- the output looks better
and scores a higher PSNR without actually helping classification.

\textbf{(ii) Super-resolve, then classify.}\quad
A special case of (i) that fixes only the resolution. As
Section~\ref{sec:tests} will show, resolution is indeed the dominant axis, so
this is the most promising restoration target; but blur and noise still cost
several \Mpp{} each at the moderate operating point, so a resolution-only
fix cannot recover the full loss on its own.

\textbf{(iii) Train-time degradation exposure (adopted).}\quad
Train the classifier directly on synthetically degraded images, so the
representation adapts to the degraded distribution. This is the approach used
by modern robustness benchmarks [6]. It needs no separate restoration
network, keeps inference to a single forward pass, and --- most important for
this study --- it exposes the architectural differences directly, because
every backbone sees exactly the same degraded distribution.

\subsection{Contributions}
This work contributes:

\begin{itemize}
  \item a 402-cell training campaign across three architectures, two
        datasets, five degradation levels, and seven experimental phases, all
        trained to convergence under a single fair-comparison invariant
        ($224 \times 224$ uniform input, frozen Optuna hyperparameters,
        identical data pipeline);
  \item a 48-replicate multi-seed audit of the 24 headline cells that bounds
        run-to-run variance below 1.2~\Mpp{};
  \item a deterministic 256-image PSNR / SSIM probe per pipeline, so that
        accuracy can be compared against measured image quality rather than
        only against the preset severity index;
  \item a per-axis attribution showing that resolution is the dominant cost
        --- roughly $4\times$ blur or salt-and-pepper at matched severity;
  \item a regularization-recovery sweep (Phase~D) and a protocol
        simplification sweep (Phase~B2 / B2-nr) that together separate the
        part of the collapse caused by destroyed information from the part
        training can recover.
\end{itemize}

\subsection{Roadmap}
Section~\ref{sec:phases} defines the seven experimental phases, which every
later section refers back to. Sections~\ref{sec:spec} and~\ref{sec:design}
give the as-built specification (with its drifts from the preliminary one)
and the engineering design. Section~\ref{sec:tests} reports the results ---
Phases B1 and C2 in full, the rest in brief, closing with the per-model
synthesis. Sections~\ref{sec:problems} and~\ref{sec:conclusions} cover the
problems encountered and the conclusions, goal evaluation, and
plan-versus-execution comparison; Section~\ref{sec:references} lists the
references.
